\subsection{Erstellung eines Kleinsignal-Ersatzschaltbilds}
% Alt: Aufgabe 21
\Aufgabe{
    \label{Aufg21}
    Zeichnen Sie das Kleinsignalersatzschaltbild der gegebenen Schaltung.
    \begin{figure}[H]
        \centering
        \input{Tikz/src/FigErsatzschaltbild2}\alttext{Die Abbildung zeigt eine elektronische Schaltung. Die elektronische Schaltung besteht aus Transistor in der Mitte. Links vom Transistor sind zwei Widerstände, \(R_2\) oben und \(R_1\) unten, die beide mit der Basis des Transistors verbunden sind. \(R_2\) und \(R_1\) wirken dabei als Spannungsteiler. Der Widerstand \(R_2\) ist außerdem mit einer positiven Spannungsquelle \(U_\mathrm{V}\) verbunden. Der Widerstand \(R_1\) ist mit Masse (Erdung) verbunden. Am linken Rand befindet sich ein Eingangspunkt \(U_\mathrm{E}\), der ebenfalls mit Masse verbunden ist. Rechts vom Transistor ist ein Ausgangspunkt \(U_\mathrm{A}\), der nach unten zur Masse führt. Der Kollektor des Transistors ist direkt mit \(U_\mathrm{V}\) verbunden, während der Emitter über den Widerstand \(R_3\) mit der Masse verbunden ist.}
        \label{fig:FigErsatzschaltbild2}
    \end{figure}
}


\Loesung{
    \begin{figure}[H]
        \centering
        \begin{subfigure}[b]{0.8\textwidth}
            \centering
            \input{Tikz/src/LsgErsatschaltbild2_1}

        \end{subfigure}\label{fig:LsgErsatschaltbild2_1}
        \vspace{0.5cm}
        \begin{subfigure}[b]{0.8\textwidth}
            \centering
            \input{Tikz/src/LsgErsatschaltbild2_2}

        \end{subfigure}\label{fig:LsgErsatschaltbild2_2}
        \label{fig:Loesung21}
    \end{figure}
    Siehe: Abschnitt \ref{fig:ErsatzschaltbildEinesBipolartransistors} und \ref{fig:ErsatzschaltbilderEmitter-Kollektor-UndBasisschaltung}.
}